% ch4_b (书页 163-174 / p178-p189)
%--A-TAIL--
% ---- 书页 163 / p178 ----（承接 4.3 节“拉格朗日表述”，前块止于书页 162 顶部式 (4.67)）
我们已经利用迹的循环性质，得以忽略 $M^{-1}$ 与 $\delta M$ 可能不对易这一事实。取矩阵 $M$ 为度规 $g_{\mu\nu}$，使得 $\det M=\det g_{\mu\nu}=g$，便得到
\begin{align*}
\delta g &= g\left(g^{\mu\nu}\delta g_{\mu\nu}\right)\\
&= -g\left(g_{\mu\nu}\delta g^{\mu\nu}\right). \tag{4.68}
\end{align*}
最后一步中，我们利用式 (4.56) 把 $\delta g_{\mu\nu}$ 化成了 $\delta g^{\mu\nu}$。现在直接代入便得到
\begin{align*}
\delta\sqrt{-g} &= -\frac{1}{2\sqrt{-g}}\,\delta g\\
&= \frac{1}{2}\,\frac{g}{\sqrt{-g}}\,g_{\mu\nu}\delta g^{\mu\nu}\\
&= -\frac{1}{2}\sqrt{-g}\,g_{\mu\nu}\delta g^{\mu\nu}. \tag{4.69}
\end{align*}
回想一下式 (4.58)，并且记得 $(\delta S)_1$ 没有贡献，我们发现
\begin{equation*}
\delta S_H = \int d^nx\sqrt{-g}\left[R_{\mu\nu}-\frac{1}{2}Rg_{\mu\nu}\right]\delta g^{\mu\nu}.
\tag{4.70}
\end{equation*}
回忆一下，作用量的泛函导数满足
\begin{equation*}
\delta S = \int\sum_i\left(\frac{\delta S}{\delta\Phi^i}\,\delta\Phi^i\right)d^nx,
\tag{4.71}
\end{equation*}
其中 $\{\Phi^i\}$ 是被变分的诸场的一个完全集合（在我们的情形下，它就只是 $g^{\mu\nu}$）。驻点（stationary points）就是让每一个 $\delta S/\delta\Phi^i=0$ 的那些点，于是我们恢复了真空中的 Einstein 方程：
\begin{equation*}
\frac{1}{\sqrt{-g}}\,\frac{\delta S_H}{\delta g^{\mu\nu}} = R_{\mu\nu}-\frac{1}{2}Rg_{\mu\nu} = 0.
\tag{4.72}
\end{equation*}

拉格朗日方法的优势体现在这样一个事实之中：我们的第一个猜测（它实际上是唯一的）就给出了正确的答案，与此前的试错法形成对照。这反映了该技术的两个优雅特征：其一，拉格朗日量是标量而非张量，因此受到的限制更强；其二，理论的对称性被直截了当地施加进来（在本例中，我们自动导出了一个散度为零的张量，这与微分同胚不变性有关，见附录 B 中的讨论）。

我们导出的 Einstein 方程是“真空中的”，因为我们在作用量里只包含了引力部分，而没有包含物质场的附加项。然而我们真正想要的，是同时得到非真空的场方程。
% ---- 书页 164 / p179 ----
这意味着我们考虑如下形式的作用量
\begin{equation*}
S = \frac{1}{16\pi G}S_H + S_M,
\tag{4.73}
\end{equation*}
其中 $S_M$ 是物质的作用量，我们极有先见之明地归一化了引力的作用量，从而恰好得到正确的答案。按照与上面相同的步骤进行下去，就得到
\begin{equation*}
\frac{1}{\sqrt{-g}}\,\frac{\delta S}{\delta g^{\mu\nu}} = \frac{1}{16\pi G}\left(R_{\mu\nu}-\frac{1}{2}Rg_{\mu\nu}\right)+\frac{1}{\sqrt{-g}}\,\frac{\delta S_M}{\delta g^{\mu\nu}} = 0.
\tag{4.74}
\end{equation*}
现在我们大胆地把能动张量（energy-momentum tensor）定义为
\begin{equation*}
T_{\mu\nu} = -2\,\frac{1}{\sqrt{-g}}\,\frac{\delta S_M}{\delta g^{\mu\nu}}.
\tag{4.75}
\end{equation*}
这让我们得以恢复完整的 Einstein 方程
\begin{equation*}
R_{\mu\nu}-\frac{1}{2}Rg_{\mu\nu} = 8\pi G T_{\mu\nu},
\tag{4.76}
\end{equation*}
或者等价地，$G_{\mu\nu}=8\pi G T_{\mu\nu}$。

我们凭什么认为式 (4.75) 真的就是能动张量呢？从某种意义上说，理由仅仅是：它是一个对称、守恒、量纲为能量密度的 (0,2) 型张量；如果你更喜欢用别的名字称呼它，请自便。但它也与我们的先验预期相吻合。再次考虑标量场的作用量，即式 (4.52)。现在对这个作用量做变分——不是对 $\phi$，而是对逆度规：
\begin{align*}
\delta S_\phi &= \int d^nx\left[\sqrt{-g}\left(-\frac{1}{2}\delta g^{\mu\nu}\nabla_\mu\phi\nabla_\nu\phi\right)+\delta\sqrt{-g}\left(-\frac{1}{2}g^{\mu\nu}\nabla_\mu\phi\nabla_\nu\phi - V(\phi)\right)\right] \tag{4.77}\\
&= \int d^nx\sqrt{-g}\,\delta g^{\mu\nu}\left[-\frac{1}{2}\nabla_\mu\phi\nabla_\nu\phi+\left(-\frac{1}{2}g_{\mu\nu}\right)\left(-\frac{1}{2}g^{\rho\sigma}\nabla_\rho\phi\nabla_\sigma\phi - V(\phi)\right)\right]. \tag{4.78}
\end{align*}
于是我们有
\begin{align*}
T^{(\phi)}_{\mu\nu} &= -2\,\frac{1}{\sqrt{-g}}\,\frac{\delta S_\phi}{\delta g^{\mu\nu}}\\
&= \nabla_\mu\phi\nabla_\nu\phi - \frac{1}{2}g_{\mu\nu}g^{\rho\sigma}\nabla_\rho\phi\nabla_\sigma\phi - g_{\mu\nu}V(\phi). \tag{4.79}
\end{align*}
在平直时空中，它恰好约化为我们曾在第 1 章中认定的、标量场的正确能动张量。

另一方面，在 Minkowski 空间中还存在能动张量的另一种定义，电磁学或场论的教科书里有时会给出它。在这种背景下，能量-动量守恒是拉格朗日量在时空平移下具有对称性的后果。\textbf{诺特定理}（Noether's theorem）断言，拉格朗日量的每一个对称性都蕴含一个守恒定律的存在；四个时空平移下的不变性导致一个张量 $S^{\mu\nu}$，它满足 $\partial_\mu S^{\mu\nu}=0$（四个关系式，$\nu$ 的每个取值对应一个）。细节可参见 Wald (1984) 或 Peskin and Schroeder (1995)。把 Noether 的步骤应用于一个依赖某些场 $\Phi^i$ 及其一阶导数 $\partial_\mu\Phi^i$ 的拉格朗日量（在平直时空中），我们得到
% ---- 书页 165 / p180 ----
\begin{equation*}
S^{\mu\nu} = \frac{\delta\mathcal{L}}{\delta(\partial_\mu\Phi^i)}\,\partial^\nu\Phi^i - \eta^{\mu\nu}\mathcal{L},
\tag{4.80}
\end{equation*}
其中对 $i$ 求和是默认的。你可以检验，凭借物质场的运动方程，这个张量是守恒的。$S^{\mu\nu}$ 常被称为“正则能动张量”（canonical energy-momentum tensor）；然而，出于若干理由，对我们来说使用式 (4.75) 更为方便。首先，当 Einstein 方程从作用量导出时，出现在方程右边的实际上正是式 (4.75)，而式 (4.80) 未必总能推广到弯曲时空。但即便在平直空间中，式 (4.75) 也有其优势：它显然是对称的，而且保证是规范不变的——这两点对式 (4.80) 都不成立。因此我们将坚持用式 (4.75) 作为能动张量的定义。

至此，Einstein 方程已经导出，本章剩下的部分将致力于探讨它的一些性质。这些讨论引人入胜，但并非严格必需；如果你愿意，可以直接跳到后续章节所讨论的应用。

\section{Einstein 方程的性质}

可以把 Einstein 方程看成度规张量场 $g_{\mu\nu}$ 的一组二阶微分方程。实际上有十个独立的方程（因为两边都是对称的两指标张量），对于度规分量的十个未知函数来说，这个数目似乎恰好合适。然而，Bianchi 恒等式 $\nabla^\mu G_{\mu\nu}=0$ 对函数 $R_{\mu\nu}(x)$ 构成四个约束，所以式 (4.44) 中只有六个真正独立的方程。事实上这正是恰当的，因为如果一个度规在某一坐标系 $x^\mu$ 中是 Einstein 方程的解，它在任何其他坐标系 $x^{\mu'}$ 中也应是其解。这意味着 $g_{\mu\nu}$ 中有四个非物理的自由度，由四个函数 $x^{\mu'}(x^\mu)$ 所代表，因此我们应当预期 Einstein 方程只约束那六个与坐标无关的自由度。

作为微分方程，它们极其复杂：Ricci 标量与 Ricci 张量是 Riemann 张量的缩并，而 Riemann 张量涉及 Christoffel 符号的导数与乘积，Christoffel 符号又涉及逆度规以及度规的导数。此外，能动张量 $T_{\mu\nu}$ 一般也会包含度规。这些方程还是非线性的，因此无法把两个已知解叠加出第三个解。所以，要想在任意一般性下求解 Einstein 方程是极其困难的，通常必须做一些简化假设。即使在令能动张量为零的真空情形，所得的方程 (4.46) 也可能很难求解。最流行的简化假设是度规具有相当程度的对称性；我们稍后会看到等距（isometries）如何让问题变得容易。

广义相对论的非线性值得一提。在 Newton 引力中，两个点质量所产生的势就是各自势的简单相加，但显然在弱场极限之外，这一点无法照搬到广义相对论中来。这有一个物理上的原因：在 GR 中，引力场与自身耦合。这可以看成等效原理的一个后果——如果引力不与自身耦合，那么一个“引力原子”（两个粒子靠彼此的引力吸引束缚在一起）的惯性质量将不同于它的引力质量（由于结合能为负）。Einstein 方程的非线性正是引力对自身之反作用（back-reaction）的一种体现。

Feynman 图为思考这一点提供了很好的方式。在量子场论中，Feynman 图被用来计算散射过程的振幅；把不同相互作用各自的贡献（每一种都由自己的图表示）加起来，就能得到这些振幅。即便我们不去量子化引力、计算散射截面（见本节末尾），也仍然可以画 Feynman 图，作为记录哪些相互作用存在、哪些不存在的一种简单方法。一个简单的例子是两个电子之间的电磁相互作用；它可以看成源于一个虚光子的交换，如图 4.1 所示。
%FIG{4.1}: {电磁相互作用的 Feynman 图。在量子场论中，这类图被用来计算散射过程的振幅；在这里，只需把它当作表示某种相互作用的示意图即可。这张图的要点在于，光子与电子的耦合正是电子之间电磁相互作用的起因。与此相反，光子与其他光子之间没有耦合，也不存在光子之间相互作用的类似图形。}

与此相反，不存在两个光子彼此之间交换另一个光子的图，因为电磁理论是线性的（没有反作用）。而引力相互作用则可以看成源于一个虚引力子（graviton，度规的量子化微扰）的交换。非线性表现为：电子与引力子都能交换虚引力子，从而产生引力，如图 4.2 所示。引力的这一特征并不独特；大多数规范理论都有它，例如作为强相互作用理论的量子色动力学（quantum chromodynamics）。电磁理论实际上是例外；其线性可以追溯到相关规范群 U(1) 是阿贝尔群这一事实。但非线性确实代表了对 Newton 理论的偏离。这一差异在实验上是可以检测的；（正如我们将看到的）水星轨道在 GR 中与在 Newton 引力中之所以不同，正是因为引力场会影响自身，而且越靠近太阳，这种影响就越明显。
% ---- 书页 167 / p182 ----
%FIG{4.2}: {引力的 Feynman 图。量子化之后，Einstein 方程预言自旋为 2 的粒子，称为引力子。我们尚不知道如何一致地完成这种量子化，但引力子的存在性足够稳固，被预期为任何良好定义的方案都会具有的特征。由于引力与能量-动量耦合，引力子与所有种类的粒子（包括其他引力子）都有相互作用。这为思考 Einstein 理论的非线性提供了一种方式。}

除了复杂与非线性之外，还有一点值得思考：Einstein 方程究竟告诉了我们什么？显然，它把能量-动量分布与曲率张量的分量联系起来；但从物理的观点看，给定某种源，究竟会产生哪一种引力场？回答这个问题的一种方式，是考虑一族相邻类时测地线的\emph{膨胀}（expansion）$\theta$ 的演化。设想一小团自由检验粒子沿测地线运动，具有四速度 $U^\mu$，我们追踪它们的演化；膨胀 $\theta=\nabla_\mu U^\mu$ 告诉我们，这团粒子的体积在任一时刻是在增长（或在 $\theta<0$ 时收缩）。显然，膨胀的数值将依赖于检验粒子的初始条件。而引力的效应则编码在膨胀的\emph{演化}之中，后者由 Raychaudhuri 方程支配。这条在附录 F 中讨论的方程告诉我们，膨胀沿测地线对固有时 $\tau$ 的导数由下面的表达式给出：
\begin{equation*}
\frac{d\theta}{d\tau} = 2\omega^2 - 2\sigma^2 - \frac{1}{3}\theta^2 - R_{\mu\nu}U^\mu U^\nu.
\tag{4.81}
\end{equation*}
右边的各项在附录 F 中有仔细的解释；$\omega$ 编码测地线的转动，$\sigma$ 编码切变，而 $R_{\mu\nu}$ 当然就是 Ricci 张量。Raychaudhuri 方程是纯几何的关系，完全不涉及 Einstein 方程。然而，把两个方程结合起来，就可以描述能量-动量如何影响检验粒子的运动，因为 Einstein 方程把 $T_{\mu\nu}$ 与 $R_{\mu\nu}$ 联系起来，而 Raychaudhuri 方程把 $R_{\mu\nu}$ 与 $d\theta/d\tau$ 联系起来。
% ---- 书页 168 / p183 ----
让我们考虑最简单的情形：一开始，所有相邻粒子在时空的一个小区域内彼此相对静止。于是在这一初始时刻，膨胀、转动和切变全都为零。我们进一步构造局部惯性坐标系 $x^{\hat\mu}$，其中 $U^{\hat\mu}$ 处于其静止系，使得 $U^{\hat\mu}=(1,0,0,0)$，且 $R_{\hat\mu\hat\nu}U^{\hat\mu}U^{\hat\nu}=R_{\hat 0\hat 0}$。于是我们（在这组坐标下、在该点）有
\begin{equation*}
\frac{d\theta}{d\tau} = -R_{\hat 0\hat 0}.
\tag{4.82}
\end{equation*}
现在我们可以求助于如下形式的 Einstein 方程：
\begin{equation*}
R_{\mu\nu} = 8\pi G\left(T_{\mu\nu}-\frac{1}{2}Tg_{\mu\nu}\right).
\tag{4.83}
\end{equation*}
由于我们处在局部惯性坐标系中，所以有
\begin{align*}
g_{\hat\mu\hat\nu} &= \eta_{\hat\mu\hat\nu} \tag{4.84}\\
T &= g^{\hat\mu\hat\nu}T_{\hat\mu\hat\nu} = -\rho + p_x + p_y + p_z, \tag{4.85}
\end{align*}
其中 $\rho=T_{\hat 0\hat 0}$ 是静止系能量密度，而 $p_k=T_{\hat k\hat k}$ 是 $x^{\hat k}$ 方向上的压强。于是，式 (4.82) 变成
\begin{equation*}
\frac{d\theta}{d\tau} = -4\pi G(\rho + p_x + p_y + p_z).
\tag{4.86}
\end{equation*}
这个方程告诉我们：能量与压强产生一种引力场，它使初始静止的检验粒子团的体积趋于缩小（如果 $\rho$ 和各个 $p_i$ 都是正的）。换言之，引力是吸引的。

当然，从式 (4.86) 我们看到，引力并不\emph{必然}是吸引的；我们可以想象 $\rho+p_x+p_y+p_z$ 为负数的源。显然，压强所扮演的角色值得注意。一方面，它代表了对 Newton 理论毫不含糊的偏离——在 Newton 理论中压强不影响引力（它不出现在 Poisson 方程 $\nabla^2\Phi=4\pi G\rho$ 中）。这一差异在我们的太阳系里很难察觉，因为太阳和行星内部的压强远小于能量密度，而后者由组成粒子的静质量主导。另一方面，请注意压强的\emph{引力}效应与我们更熟悉的\emph{直接}效应正好相反——直接效应是正压强会把东西推开。在多数情形下，压强的直接效应要明显得多。然而，压强只有在存在压强梯度时才能直接起作用（例如活塞内外之间的压强差），而引力效应只依赖压强在当地的数值。如果压强完全平滑，它就只能通过其引力效应被探测到；真空能就是这样的一个例子，我们将在 4.5 节讨论它。
% ---- 书页 169 / p184 ----
关于式 (4.86) 最后再评论一点：它与 Einstein 方程完全等价——两者传达完全相同的信息。这个非常具体的关系对任何一组初始静止的检验粒子都成立；而它能成立的唯一途径，就是 Einstein 方程的所有分量都为真。因此，如果我们愿意，就可以用文字\footnote{J.C. Baez, “The Meaning of Einstein's Equation”〔“Einstein 方程的含义”〕， http://arXiv.org/abs/gr-qc/0103044.}把 Einstein 方程表述如下：“任何一组初始静止的粒子，其体积的膨胀正比于（负的）能量密度与三个压强分量之和。”

于是，Einstein 方程告诉我们，能量密度和压强以某种方式影响 Ricci 张量，使得当 $\rho$ 和 $p$ 为正时粒子相互吸引。那么，Riemann 张量中不包含在 Ricci 张量之内的那些分量呢？在第 3 章中我们发现，这些分量由 Weyl 张量描述（这里写出四维情形的表达式）：
\begin{equation*}
C_{\rho\sigma\mu\nu} = R_{\rho\sigma\mu\nu} + \frac{1}{3}g_{\rho[\mu}g_{\nu]\sigma}R - g_{\rho[\mu}R_{\nu]\sigma} + g_{\sigma[\mu}R_{\nu]\rho}.
\tag{4.87}
\end{equation*}
Ricci 张量是 Riemann 张量的迹，而 Weyl 张量描述其无迹部分；二者合起来完整地刻画了曲率。显然，给定某个确定的能量-动量分布，Weyl 曲率的选择仍有某种自由，因为没有类似 Einstein 方程的关系能把 $C^\rho{}_{\sigma\mu\nu}$ 与 $T_{\mu\nu}$ 代数地联系起来。而这正是理应如此。例如，设想一个处处真空的时空，$R_{\mu\nu}=0$。平直的 Minkowski 空间是这种情形下的一个可能解，但穿越空旷时空传播的引力波同样是可能的解（我们将在第 7 章讨论）。

既然进入 Einstein 方程的只有 $R_{\mu\nu}$，看起来 $C_{\rho\sigma\mu\nu}$ 的分量似乎完全不受约束。但请回忆，我们并不被允许在流形上任意指定曲率张量的分量；它们由 Bianchi 恒等式联系着，
\begin{equation*}
\nabla_{[\lambda}R_{\rho\sigma]\mu\nu} = 0.
\tag{4.88}
\end{equation*}
正如你在第 3 章习题 10 中证明过的，这个恒等式蕴含 Weyl 张量满足如下形式的微分关系：
\begin{equation*}
\nabla^\rho C_{\rho\sigma\mu\nu} = \nabla_{[\mu}R_{\nu]\sigma} + \frac{1}{6}g_{\sigma[\mu}\nabla_{\nu]}R.
\tag{4.89}
\end{equation*}
在右边，Riemann 张量只是通过它的缩并——Ricci 标量与 Ricci 张量——出现，而后者可以经由 Einstein 方程与 $T_{\mu\nu}$ 联系起来；于是我们有
\begin{equation*}
\nabla^\rho C_{\rho\sigma\mu\nu} = 8\pi G\left(\nabla_{[\mu}T_{\nu]\sigma} + \frac{1}{3}g_{\sigma[\mu}\nabla_{\nu]}T\right).
\tag{4.90}
\end{equation*}
因此，虽然 $R_{\mu\nu}$ 与 $T_{\mu\nu}$ 通过 Einstein 方程代数地相联系，$C_{\rho\sigma\mu\nu}$ 与 $T_{\mu\nu}$ 却是由这个一阶微分方程联系的。对一个给定的能量-动量分布，会存在若干可能的解，每个解由特定的边界条件确定。这个方程可以看成引力波的传播方程，与麦克斯韦方程组 $\nabla_\mu F^{\nu\mu}=J^\nu$ 极为类似。
% ---- 书页 170 / p185 ----
列举了 Einstein 方程这么多可爱的性质之后，似乎也应当公平地提一提它令人沮丧的一面：广为人知的、广义相对论与量子力学难以调和的困难。GR 是一个经典场论：动力学变量是定义在时空上的一个场（度规），由这个场构造出的坐标不变量（如标量曲率）原则上可以被任意精确地指定和测量。对于其他场论（如电磁学），从经典理论出发将其量子化、进而得到作用在 Hilbert 空间中的波函数之上的算符的动力学，有着理解透彻的成熟程序。而对 GR 来说，通常的程序会同时遭遇技术困难和概念困难，对此的描述超出了本书的范围。技术困难的一个方面是，GR 不像粒子物理的标准模型那样是“可重整化的”（renormalizable）；当考虑高阶量子效应时，会出现无法被任何有限多个参数吸收的无穷大。不可重整化并不意味着理论在根本上就是错的，但它强烈暗示，只应在某个能标以下才认真地对待这个理论。

幸运的是，量子引力的可观测效应预计要远离我们的日常经验（或者顺便说一句，也远离实验室中能实现的任何条件）才会变得重要。早在 1899 年，Planck 就注意到，他自己的常数 $h$——如今我们更常用 $\hbar = h/2\pi = 1.05\times10^{-27}\,\mathrm{cm}^2\,\mathrm{g/sec}$ 来代替它——可以与 Newton 引力常数 $G = 6.67\times10^{-8}\,\mathrm{cm}^3\,\mathrm{g}^{-1}\,\mathrm{sec}^{-2}$ 以及光速 $c = 3.00\times10^{10}\,\mathrm{cm}\,\mathrm{sec}^{-1}$ 组合起来，构成一组基本的带量纲的量：Planck 质量
\begin{equation*}
m_{\mathrm{P}} = \left(\frac{\hbar c}{G}\right)^{1/2} = 2.18\times 10^{-5}\,\mathrm{g},
\tag{4.91}
\end{equation*}
Planck 长度
\begin{equation*}
l_{\mathrm{P}} = \left(\frac{\hbar G}{c^3}\right)^{1/2} = 1.62\times 10^{-33}\,\mathrm{cm},
\tag{4.92}
\end{equation*}
Planck 时间
\begin{equation*}
t_{\mathrm{P}} = \left(\frac{\hbar G}{c^5}\right)^{1/2} = 5.39\times 10^{-44}\,\mathrm{sec},
\tag{4.93}
\end{equation*}
以及 Planck 能量
\begin{align*}
E_{\mathrm{P}} &= \left(\frac{\hbar c^5}{G}\right)^{1/2} = 1.95\times 10^{16}\,\mathrm{erg} \tag{4.94}\\
&= 1.22\times 10^{19}\,\mathrm{GeV}. \tag{4.95}
\end{align*}
% ---- 书页 171 / p186 ----
一个 GeV 是 $10^9$ 电子伏特，它是粒子物理中的常用单位，因为它大约是一个质子的质量。我们通常取 $\hbar = c = 1$，于是这些量就变得无从区分，即 $m_{\mathrm{P}} = l_{\mathrm{P}}^{-1} = t_{\mathrm{P}}^{-1} = E_{\mathrm{P}}$。你会听到人们说诸如“Planck 质量是 $10^{19}$ GeV”之类的话，或者干脆说“Planck 标度”（Planck scale）。另一个常用量是约化 Planck 标度（reduced Planck scale）$\bar{m}_{\mathrm{P}} = m_{\mathrm{P}}/\sqrt{8\pi} = 2.43\times10^{18}\,\mathrm{GeV}$，它在方程中往往更加方便——注意式 (4.73) 中标量曲率的系数正是 $\bar{m}_{\mathrm{P}}^2/2$。很可能，只有当我们考虑的粒子质量大于 $m_{\mathrm{P}}$、时间短于 $t_{\mathrm{P}}$、长度小于 $l_{\mathrm{P}}$、或能量高于 $E_{\mathrm{P}}$ 时，量子引力才变得重要；在更低的标度下，经典 GR 就应当足够了。由于这些都远离可观测的现象，构造一个自洽的量子引力理论更多是一个原则问题，而非实际问题。另一方面，弯曲时空中的量子效应在现实世界中可能很重要；正如我们将在第 8 章讨论的，它们可能导致早期宇宙中的密度涨落，而这些涨落会成长为今天我们观测到的星系和大尺度结构。

在完全量子的理论中，有一个在适当极限下能涵盖 GR 的领先候选者：弦论（string theory）。在弦论中，我们设想基本的对象不是像电子或光子那样的点粒子，而是被称为弦的细小的一维对象，它们既可以是闭合的圈，也可以是开放的线段。弦论最初是作为强核力的模型被提出的，但人们很快意识到，该理论不可避免地预言了一个无质量的自旋 2 粒子：这恰恰是一个量子引力理论所需要的东西。弦论看起来是一个自洽的量子理论，而且它预言了引力，但关于它仍有许多我们不理解的地方。特别是，经典时空如何从基本的弦之中产生，这一点颇为神秘；而它与直接实验的联系，充其量也十分微弱。尽管如此，弦论异常丰富而强健，有望在可预见的未来成为理论物理的重要部分。

\section{宇宙学常数}

广义相对论的一个特征是，引力场的源是整个能动张量。在非引力的物理中，只有能量从一个态到另一个态的\emph{改变}才是可测的；能量的归一化是任意的。例如，势能为 $V(x)$ 的粒子的运动，与势能为 $V(x)+V_0$ 的粒子的运动完全相同，其中 $V_0$ 是任意常数。然而在引力中，能量的实际数值是重要的，而不仅仅是态之间的差。

这种行为开启了\textbf{真空能}（vacuum energy）的可能性：一种为空的空间所特有的能量密度。我们或许希望真空具备的一个特征是：它不挑出任何一个特殊的方向；只要相应的能动张量在局部惯性坐标系中是洛伦兹不变的，非零的能量密度就仍然是可能的。洛伦兹不变性意味着相应的能动张量应当正比于度规：
% ---- 书页 172 / p187 ----
\begin{equation*}
T^{(\mathrm{vac})}_{\hat\mu\hat\nu} = -\rho_{\mathrm{vac}}\,\eta_{\hat\mu\hat\nu},
\tag{4.96}
\end{equation*}
因为 $\eta_{\hat\mu\hat\nu}$ 是唯一的洛伦兹不变 (0,2) 型张量。这可以直截了当地从惯性坐标推广到任意坐标：
\begin{equation*}
T^{(\mathrm{vac})}_{\mu\nu} = -\rho_{\mathrm{vac}}\,g_{\mu\nu}.
\tag{4.97}
\end{equation*}
与理想流体的能动张量 $T_{\mu\nu}=(\rho+p)U_\mu U_\nu + pg_{\mu\nu}$ 相比较，我们发现真空看起来像一种理想流体，其各向同性的压强与能量密度符号相反：
\begin{equation*}
p_{\mathrm{vac}} = -\rho_{\mathrm{vac}}.
\tag{4.98}
\end{equation*}
能量密度应当在整个时空中为常数，因为若有梯度就不再是洛伦兹不变的了。

如果我们把能动张量分解为物质部分 $T^{(\mathrm{M})}_{\mu\nu}$ 与真空部分 $T^{(\mathrm{vac})}_{\mu\nu}=-\rho_{\mathrm{vac}}g_{\mu\nu}$，Einstein 方程就是
\begin{equation*}
R_{\mu\nu}-\frac{1}{2}Rg_{\mu\nu} = 8\pi G\left(T^{(\mathrm{M})}_{\mu\nu}-\rho_{\mathrm{vac}}g_{\mu\nu}\right).
\tag{4.99}
\end{equation*}
在发明 GR 之后不久，Einstein 试图寻找一个静态的宇宙学模型，因为当时的天文观测似乎正暗示这一点。其结果就是 Einstein 静态宇宙，我们将在第 8 章讨论它。为了让这个静态宇宙学在普通物质源之下求解场方程，有必要添加一个称为\textbf{宇宙学常数}（cosmological constant）的新项 $\Lambda$，它以如下方式进入：
\begin{equation*}
R_{\mu\nu}-\frac{1}{2}Rg_{\mu\nu} + \Lambda g_{\mu\nu} = 8\pi G T_{\mu\nu}.
\tag{4.100}
\end{equation*}
与式 (4.99) 比较可见，宇宙学常数恰恰等价于引入一个真空能密度
\begin{equation*}
\rho_{\mathrm{vac}} = \frac{\Lambda}{8\pi G}.
\tag{4.101}
\end{equation*}
“宇宙学常数”与“真空能”这两个术语实质上可以互换。

非零的真空能是我们应当预期的东西吗？我们当初是通过寻找能从度规构造出的最简单的标量，而得到 Hilbert 拉格朗日量 $\widehat{\mathcal{L}}_H = R$ 的。当然，还有一个更简单的候选，即常数。利用式 (4.69)，不难检验
\begin{equation*}
S = \int d^4x\sqrt{-g}\left[\frac{1}{16\pi G}(R-2\Lambda) + \widehat{\mathcal{L}}_{\mathrm{M}}\right]
\tag{4.102}
\end{equation*}
% ---- 书页 173 / p188 ----
会导致修改后的方程 (4.100)；或者换个说法，真空的拉格朗日量就是
\begin{equation*}
\widehat{\mathcal{L}}_{\mathrm{vac}} = -\rho_{\mathrm{vac}}.
\tag{4.103}
\end{equation*}
所以，引入真空能当然容易；然而，对它的预期取值我们毫无洞察，因为它以一个任意常数的形式进入。

真空能归根结底本身就是自然界的一个常数。（在某些理论中会出现例外：某个时空对称性——例如超对称性或共形不变性——支配着真空能的取值；这里我们考虑的是更一般的场论。）尽管如此，真空能有各种各样不同的贡献来源，如果总取值远小于各个单独的贡献，那就奇怪了。其中一种贡献来自零点涨落（zero-point fluctuations）——量子场处于真空态时的能量。

考虑一个简谐振子，即在一维势 $V(x)=\frac{1}{2}\omega^2 x^2$ 中运动的一个粒子。经典地看，这个系统的真空是粒子静止于势的极小值处（$x=0$）的态，此时能量为零。然而从量子力学上看，不确定性原理不允许我们同时确定粒子的位置和动量，我们会发现最低能量态具有能量 $E_0=\frac{1}{2}\hbar\omega$（这里为了清晰起见，我们临时重新引入了显式的 $\hbar$ 因子）。当然，在没有引力时，无论哪个系统，其真空能实际上都是完全任意的；我们可以给势加上任何常数而不改变理论。但量子涨落已经使零点能不同于我们的经典预期。

在场论中有一个完全类似的情形。如果我们对一个自由量子场（为简单起见忽略相互作用的场）做傅里叶变换，就会发现它变成动量空间中无穷多个谐振子，正如我们将在第 9 章讨论的。每个振子的频率为 $\omega=\sqrt{m^2+k^2}$，其中 $m$ 是场的质量，$k$ 是该模式波矢的大小。如果我们把经典真空能取为零，这些模式中的每一个都贡献一份量子零点能 $\hbar\omega/2$。形式上，把所有这些贡献加起来会得到一个无穷大的结果。然而，如果我们以“只在某个紫外动量截断 $k_{\mathrm{max}}$ 以下才信任我们的理论”为由，丢弃动量非常高的那些模式，就会发现所得的能量密度具有如下形式：
\begin{equation*}
\rho_{\mathrm{vac}} \sim \hbar k_{\mathrm{max}}^4.
\tag{4.104}
\end{equation*}
这个答案本来可以靠量纲分析猜出来；被略去的数值常数将依赖于所考虑的具体理论。如果我们有信心把普通量子场论一直用到约化 Planck 标度 $\bar{m}_{\mathrm{P}} = (8\pi G)^{-1/2} \sim 10^{18}\,\mathrm{GeV}$，我们预计会有一个量级为
\begin{equation*}
\rho_{\mathrm{vac}} \sim (10^{18}\,\mathrm{GeV})^4 \sim 10^{112}\,\mathrm{erg/cm^3}.
\tag{4.105}
\end{equation*}
% ---- 书页 174 / p189 ----
的贡献。场论也可能更早就失效，不过，要相信它会在任何具体的标度上失效，量子引力是我们所拥有的最好的理由。

正如我们将在第 8 章讨论的，宇宙学观测意味着
\begin{equation*}
\left|\rho^{(\mathrm{obs})}_{\Lambda}\right| \leq (10^{-12}\,\mathrm{GeV})^4 \sim 10^{-8}\,\mathrm{erg/cm^3},
\tag{4.106}
\end{equation*}
这比刚才导出的朴素预期小得多。式 (4.105) 与式 (4.106) 之比，正是宇宙学常数的理论值与观测值之间那著名的 120 个数量级差异的由来。我们可以随意设想裸真空能经过调整，使得净宇宙学常数与极限 (4.106) 相符；但有一个问题：我们不知道有任何特殊的对称性，既能强制真空能为零，又与已知的物理定律相容；这个难题就是“宇宙学常数问题”（cosmological constant problem）。我们将在第 8 章更多地讨论真空能的宇宙学效应。\footnote{关于真空能的物理与宇宙学的更多内容，见 S.M. Carroll, \emph{Liv. Rev. Rel.} \textbf{4}, 1 (2001), http://arxiv.org/astro-ph/0004075.}

\section{能量条件}

有时，不指明 $T_{\mu\nu}$ 由哪一种物质理论导出，而去思考 Einstein 方程，是有用的。这会给我们留下极大的任意性；例如考虑这样一个问题：哪些度规服从 Einstein 方程？在没有对 $T_{\mu\nu}$ 的某些约束时，答案是任何度规都行；只需取你喜欢的度规，为它算出 Einstein 张量 $G_{\mu\nu}$，然后要求 $T_{\mu\nu}$ 等于 $G_{\mu\nu}$ 即可。而由 Bianchi 恒等式，它自动守恒。我们真正关心的，是在“现实的”（realistic）能量与动量源——姑且不论这个词是什么意思——存在时 Einstein 方程解的存在性。一种策略是考虑具体的源，比如标量场、尘埃或电磁场。然而，我们有时希望理解的是那些对各种不同源都成立的 Einstein 方程的性质。在这种情况下，方便的做法是施加\emph{能量条件}（energy conditions）来限制 $T_{\mu\nu}$ 的任意性。

能量条件是对能动张量的坐标不变的限制。因此我们必须从 $T_{\mu\nu}$ 构造标量，这通常通过把它与任意的类时矢量 $t^\mu$ 或类光矢量 $\ell^\mu$ 缩并来实现。例如，弱能量条件（weak energy condition, WEC）陈述：对所有类时矢量 $t^\mu$，都有 $T_{\mu\nu}t^\mu t^\nu\geq 0$。为了建立物理直觉，考虑源为理想流体的特殊情形是有用的，此时能动张量取如下形式：
\begin{equation*}
T_{\mu\nu} = (\rho + p)U_\mu U_\nu + pg_{\mu\nu},
\tag{4.107}
\end{equation*}
其中 $U^\mu$ 是流体的四速度。让我们用这个形式把 WEC 翻译成物理的语言。由于压强是各向同性的，只要对某个类光矢量 $\ell^\mu$，$T_{\mu\nu}U^\mu U^\nu \geq 0$ 与 $T_{\mu\nu}\ell^\mu\ell^\nu \geq 0$ 同时成立，
% 本块末句在原书书页174底部中断（“will be nonnegative”续接书页175顶部 “for all timelike vectors t^mu if both T_{mu nu}U^mu U^nu >= 0 and ...”），下一块请用 %JOIN% 续接本句。
